\documentclass[11pt]{article}
\usepackage{metricsstyle}

%% Beginner cheat sheet for ECON*6140 Assignment #1.
%% Every regression is done by matrix algebra (numpy only, no statsmodels).
%% All snippets were run against tbrate.xlsx in "Metrics HW1" before typesetting.
%% The previous, longer statsmodels version is pythonguide_long.tex.
%% This file is NOT part of the lesson/hw build glob in build.sh.

\usepackage{listings}
\usepackage{xcolor}
\usepackage{booktabs}
\usepackage{array}
\usepackage[hidelinks]{hyperref}
\usepackage{microtype}
\usepackage{upquote}
\setlength{\emergencystretch}{3.2em}

\definecolor{codebg}{RGB}{246,246,248}
\definecolor{codekw}{RGB}{0,83,166}
\definecolor{codestr}{RGB}{163,21,21}
\definecolor{codecom}{RGB}{0,120,60}
\definecolor{codegray}{RGB}{120,120,120}

\lstdefinestyle{py}{
  language=Python,
  backgroundcolor=\color{codebg},
  basicstyle=\ttfamily\footnotesize,
  keywordstyle=\color{codekw}\bfseries,
  stringstyle=\color{codestr},
  commentstyle=\color{codecom}\itshape,
  frame=single,
  framerule=0.3pt,
  rulecolor=\color{codegray!60},
  breaklines=true,
  showstringspaces=false,
  columns=fullflexible,
  keepspaces=true,
  upquote=true,
  aboveskip=6pt,
  belowskip=6pt,
}
\lstset{style=py}

\newcommand{\code}[1]{\texttt{\detokenize{#1}}}
\setlength{\tabcolsep}{5pt}
\renewcommand{\arraystretch}{1.2}

\begin{document}

\begin{center}
  {\large\scshape ECON*6140 Econometrics I}\\[3pt]
  {\LARGE\bfseries Python Cheat Sheet for HW1}\\[4pt]
  {\itshape Matrices only --- no regression packages}
\end{center}
\vspace{0.3em}\hrule\vspace{0.6em}

\noindent\textbf{How to use it.} Copy one grey block into one notebook cell,
change the names, press \textbf{Shift+Enter}. One idea per cell, so when something
breaks you know where. Run cells top to bottom.

%--------------------------------------------------------------------
\section{Five rules}

\begin{flatnum}
  \item \code{@} is matrix multiplication. \code{*} multiplies element by element.
        Mixing them up gives a wrong answer with \emph{no} error.
  \item \code{X.T} is $X'$. \code{np.linalg.inv(A)} is $A^{-1}$. \code{np.eye(n)} is the
        $n\times n$ identity $I$.
  \item Python counts from 0. \code{b[0]} is the first coefficient (the constant),
        \code{b[1]} the second.
  \item \code{df['r']} picks a column (square brackets). \code{df.head()} runs an
        action (round brackets).
  \item \code{X.shape} tells you the size of a matrix. Print it whenever you are unsure.
\end{flatnum}

%--------------------------------------------------------------------
\section{Load the data and make variables}

\begin{lstlisting}
import numpy as np
import pandas as pd
import matplotlib.pyplot as plt
np.set_printoptions(suppress=True)      # show 0.0161, not 1.61e-02

df = pd.read_excel('tbrate.xlsx', sheet_name='data')
df.index = pd.PeriodIndex(df['period'], freq='Q')   # rows named 1950Q1, 1950Q2, ...
\end{lstlisting}

\noindent Two tools make every variable in the assignment:

\begin{center}
\begin{tabular}{@{}lll@{}}
\toprule
Code & Means & Example \\
\midrule
\code{.diff()}   & change from last period, $\Delta x_t = x_t - x_{t-1}$ & \code{df['r'].diff()} \\
\code{.shift(1)} & last period's value, $x_{t-1}$ & \code{df['pi'].shift(1)} \\
\code{.shift(2)} & two periods back, $x_{t-2}$ & \code{df['dr'].shift(2)} \\
\code{np.log()}  & natural log & \code{np.log(df['c'])} \\
\bottomrule
\end{tabular}
\end{center}

\begin{lstlisting}
df['dr']   = df['r'].diff()       # change in r
df['dy']   = df['y'].diff()       # change in y
df['pi_1'] = df['pi'].shift(1)    # last quarter's pi
df['dy_1'] = df['dy'].shift(1)    # last quarter's dy (needs dy made first!)
df['dr_1'] = df['dr'].shift(1)    # last quarter's dr
df['dr_2'] = df['dr'].shift(2)    # dr two quarters back
\end{lstlisting}

\noindent Lags leave empty cells (\code{NaN}) at the top. That is why the sample
starts later. Cut the sample by date (both ends included):

\begin{lstlisting}
sub = df.loc['1950Q4':'1996Q4']
len(sub)                          # 185 rows
sub.isna().sum()                  # every column should say 0
\end{lstlisting}

%--------------------------------------------------------------------
\section{Regression with matrices}

\topic{Step 1: turn columns into $y$ and $X$.} The constant is just a column of ones.

\begin{lstlisting}
one = np.ones(len(sub))                      # the constant
y = sub['dr'].to_numpy()                     # y as a plain vector
X = np.column_stack([one, sub['pi_1'], sub['dy_1'], sub['dr_1'], sub['dr_2']])
X.shape                                      # (185, 5)
\end{lstlisting}

\noindent \code{np.column_stack([...])} glues columns side by side. The order you
list them is the order of $b$.

\topic{Step 2: the formulas.}

\begin{center}
\begin{tabular}{@{}lll@{}}
\toprule
Math & Code & Name \\
\midrule
$b = (X'X)^{-1}X'y$           & \code{b = np.linalg.inv(X.T @ X) @ X.T @ y} & coefficients \\
$\hat y = Xb$                 & \code{yhat = X @ b}                         & fitted values \\
$e = y - \hat y$              & \code{e = y - yhat}                         & residuals \\
$s^2 = e'e/(n-k)$             & \code{s2 = (e @ e) / (n - k)}               & error variance \\
$\widehat{\Var}(b) = s^2(X'X)^{-1}$ & \code{V = s2 * np.linalg.inv(X.T @ X)} & covariance matrix \\
$\text{se}(b_j)=\sqrt{V_{jj}}$ & \code{se = np.sqrt(np.diag(V))}            & standard errors \\
\bottomrule
\end{tabular}
\end{center}

\topic{Step 3: write it once, reuse it.} Put this in one cell and run it. After that,
every regression in the assignment is one line.

\begin{lstlisting}
def ols(y, X):
    XtX_inv = np.linalg.inv(X.T @ X)
    b    = XtX_inv @ X.T @ y          # coefficients
    yhat = X @ b                      # fitted values
    e    = y - yhat                   # residuals
    n, k = X.shape                    # rows, columns
    s2   = (e @ e) / (n - k)
    V    = s2 * XtX_inv               # Var(b)
    se   = np.sqrt(np.diag(V))        # standard errors
    return b, se, yhat, e, V
\end{lstlisting}

\begin{lstlisting}
b, se, yhat, e, V = ols(y, X)
print(b)                          # [const, pi_1, dy_1, dr_1, dr_2]
print(se)
\end{lstlisting}

\noindent The t-statistic is \code{b / se}. You can leave pieces you don't need
with \code{_}, e.g.\ \code{_, _, _, e2, _ = ols(y2, X2)} keeps only the residuals.

\topic{Regressing one vector on another.} Anything you computed (residuals, fitted
values) is just a vector, so it goes straight into \code{column_stack}:

\begin{lstlisting}
Z = np.column_stack([one, yhat])   # a constant and the fitted values
b_z, se_z, _, _, _ = ols(e, Z)     # regress e on them
\end{lstlisting}

\noindent \textbf{No constant?} Leave out \code{one}. With a single regressor, make it
a one-column matrix: \code{X = v.reshape(-1, 1)}.

%--------------------------------------------------------------------
\section{Plot against time}

\begin{lstlisting}
t = sub.index.to_timestamp()                   # dates for the x-axis

fig, ax = plt.subplots(2, 1, figsize=(9, 6), sharex=True)
ax[0].plot(t, yhat);  ax[0].set_ylabel('fitted')
ax[1].plot(t, e);     ax[1].set_ylabel('residual')
ax[1].axhline(0, color='grey')                 # zero line
plt.show()
\end{lstlisting}

\noindent Save it for your write-up with \code{fig.savefig('q2.png', dpi=200)}.

%--------------------------------------------------------------------
\section{Projection matrices $P$ and $M$}

\begin{lstlisting}
X1 = np.column_stack([one, sub['dy_1']])       # any group of columns
P1 = X1 @ np.linalg.inv(X1.T @ X1) @ X1.T      # P1 = X1 (X1'X1)^-1 X1'
M1 = np.eye(len(sub)) - P1                     # M1 = I - P1
\end{lstlisting}

\noindent These are $n\times n$ ($185\times185$). Use them just like the math:
\code{P1 @ y}, \code{M1 @ y}, \code{M1 @ X2}. Handy facts to check your code:
$P_1 y$ is the fitted value and $M_1 y$ the residual from regressing $y$ on $X_1$.

%--------------------------------------------------------------------
\section{A combination of coefficients: $\hat\gamma = w'b$}

Make $w$ the same length as $b$, with a 1 on each coefficient you want to add up.

\begin{lstlisting}
w = np.zeros(len(b))        # [0, 0, 0, 0, 0]
w[1] = 1                    # pick b[1]
w[2] = 1                    # pick b[2]      -> w = [0, 1, 1, 0, 0]
                            # (w[1:4] = 1 would pick b[1], b[2], b[3])

gamma    = w @ b                    # w'b
se_gamma = np.sqrt(w @ V @ w)       # sqrt(w' Var(b) w)
\end{lstlisting}

\topic{Reparameterising.} To get the same number straight out of a regression,
build new columns (e.g.\ \code{sub['x'] - sub['z']}), put them in a new $X$, and run
\code{ols} again. The fitted values must not change --- use that as your check.

%--------------------------------------------------------------------
\section{Checking your answers}

\begin{lstlisting}
np.allclose(a, c)          # True if two vectors/matrices are equal (up to rounding)
np.allclose(P1 @ P1, P1)   # a projection times itself is itself -> True
np.allclose(M1 @ X1, 0)    # M1 wipes out X1 -> True
e.mean()                   # about 0 when X has a constant
\end{lstlisting}

\noindent Use \code{np.allclose}, not \code{==}: computers round, so two ``equal''
numbers can differ in the 15th decimal.

%--------------------------------------------------------------------
\section{Errors you will hit}

\begin{center}
\begin{tabular}{@{}p{0.36\linewidth}p{0.58\linewidth}@{}}
\toprule
Message & Fix \\
\midrule
\code{KeyError: 'dr'} & That column doesn't exist yet. Run the cells from the top
  (Kernel $\to$ Restart \& Run All). \\
\code{matmul: ... mismatch} & Sizes don't fit. Print \code{X.shape} and \code{y.shape}. \\
\code{Singular matrix} & Two columns are the same thing (e.g.\ \code{one} added twice). \\
\code{nan} in your answer & A \code{NaN} slipped into the sample. Check
  \code{sub.isna().sum()}. \\
\code{NameError: name 'np'} & Run the import cell first. \\
Nothing is shown & Only the \emph{last} line of a cell is displayed. Use \code{print(...)}. \\
\bottomrule
\end{tabular}
\end{center}

\end{document}
